% GENERATED FILE. Edit canonical lessons, then run npm run generate:books.
% canonical-source-hash: fnv1a64:1a8dd19fd647fb37
% canonical-lessons: PA-C324-bamh, PA-C324-modha, PA-C324-haphna, PA-C324-pasina, PA-C324-chakkar

\chapter{Aching and Out of Breath}
\label{ch:pa-a2-aching-and-out-of-breath}
\hlchaptermodality{324}

\begin{chapteropening}
\textbf{By the end of this chapter:} \emph{I can say that I am out of breath and where it hurts.}

Complete the last lesson of chapter 324: I can say that I am out of breath and where it hurts.
\end{chapteropening}

\section[\texorpdfstring{\pa{ਬਾਂਹ} (bāṁh)}{ਬਾਂਹ (bāṁh)}]{\pa{ਬਾਂਹ} (bāṁh) — an arm}

\label{lesson:PA-C324-bamh}



\begin{quote}
Before the new one: say the Punjabi for a rib, then the Punjabi for to ache.
\end{quote}

\subsection*{You'll want to know: \pa{ਬਾਂਹ}}
\textbf{\pa{ਬਾਂਹ}} — \emph{bāṁh} — \textquotedblleft{}an arm\textquotedblright{}.

\begin{quote}
\pa{ਭਾਰੀ} \pa{ਥੈਲੇ} \pa{ਤੋਂ} \pa{ਮੇਰੀ} \pa{ਬਾਂਹ} \pa{ਦੁਖਦੀ} \pa{ਹੈ।}
\end{quote}

\begin{quote}
\emph{bhārī thaile ton merī bāṁh dukhdī hai.}
\end{quote}

\begin{quote}
My arm aches from the heavy bag.
\end{quote}

A feminine noun.

\begin{cousinweb}
\textbf{\pa{ਬਾਂਹ}} (\emph{bāṁh}) is inherited from Sanskrit \emph{bāhu}, \textquotedblleft{}an arm\textquotedblright{}; CDIAL lists Punjabi \emph{bā̃h}, \textquotedblleft{}arm\textquotedblright{}, under it (CDIAL 9229).
\end{cousinweb}

\begin{grammarlens}[title={What you've built}]
One more word for this chapter.
\end{grammarlens}

\subsection*{Your turn}
\begin{itemize}
\raggedright
  \item \emph{Say it:} \emph{bāṁh}
  \item \emph{Say it:} \emph{bāṁh}, once more
  \item \emph{Say it:} \emph{bhārī thaile ton merī bāṁh dukhdī hai.} — My arm aches from the heavy bag.
  \item \emph{Recall:} say the Punjabi for a rib, then the Punjabi for to ache, then say \emph{bāṁh} again
\end{itemize}

\begin{tcolorbox}[breakable,colback=teal!4,colframe=teal!35!black,title={Before you move on}]
What does \pa{ਬਾਂਹ} mean? (\textquotedblleft{}An arm\textquotedblright{}.) Say it once more.
\end{tcolorbox}
\section[\texorpdfstring{\pa{ਮੋਢਾ} (moḍhā)}{ਮੋਢਾ (moḍhā)}]{\pa{ਮੋਢਾ} (moḍhā) — a shoulder}

\label{lesson:PA-C324-modha}



\begin{quote}
Before the new one: say the Punjabi for to ache, then the Punjabi for an arm.
\end{quote}

\subsection*{You'll want to know: \pa{ਮੋਢਾ}}
\textbf{\pa{ਮੋਢਾ}} — \emph{moḍhā} — \textquotedblleft{}a shoulder\textquotedblright{}.

\begin{quote}
\pa{ਪਿਤਾ} \pa{ਜੀ} \pa{ਬੱਚੇ} \pa{ਨੂੰ} \pa{ਮੋਢੇ} \pa{ਉੱਤੇ} \pa{ਚੁੱਕਦੇ} \pa{ਹਨ।}
\end{quote}

\begin{quote}
\emph{pitā jī bachche nū̃ moḍhe utte chukkde han.}
\end{quote}

\begin{quote}
Father carries the child on his shoulder.
\end{quote}

A masculine noun.

\begin{cousinweb}
\textbf{\pa{ਮੋਢਾ}} (\emph{moḍhā}) goes back to a reconstructed \emph{*maurdhaka}, \textquotedblleft{}pertaining to the head\textquotedblright{}; CDIAL gives no Punjabi form, but lists Lahnda \emph{moḍhā}, \textquotedblleft{}shoulder\textquotedblright{} (CDIAL 10372).
\end{cousinweb}

\begin{grammarlens}[title={What you've built}]
One more word for this chapter.
\end{grammarlens}

\subsection*{Your turn}
\begin{itemize}
\raggedright
  \item \emph{Say it:} \emph{moḍhā}
  \item \emph{Say it:} \emph{moḍhā}, once more
  \item \emph{Say it:} \emph{pitā jī bachche nū̃ moḍhe utte chukkde han.} — Father carries the child on his shoulder.
  \item \emph{Recall:} say the Punjabi for to ache, then the Punjabi for an arm, then say \emph{moḍhā} again
  \item \emph{Write it:} copy \textbf{\pa{ਮੋਢਾ}} once, checking the whole word
\end{itemize}

\begin{tcolorbox}[breakable,colback=teal!4,colframe=teal!35!black,title={Before you move on}]
What does \pa{ਮੋਢਾ} mean? (\textquotedblleft{}A shoulder\textquotedblright{}.) Say it once more.
\end{tcolorbox}
\section[\texorpdfstring{\pa{ਹਫਣਾ} (haphṇā)}{ਹਫਣਾ (haphṇā)}]{\pa{ਹਫਣਾ} (haphṇā) — to pant}

\label{lesson:PA-C324-haphna}



\begin{quote}
Before the new one: say the Punjabi for an arm, then the Punjabi for a shoulder.
\end{quote}

\subsection*{You'll want to know: \pa{ਹਫਣਾ}}
\textbf{\pa{ਹਫਣਾ}} — \emph{haphṇā} — \textquotedblleft{}to pant\textquotedblright{}.

\begin{quote}
\pa{ਦਾਦਾ} \pa{ਜੀ} \pa{ਪੌੜੀ} \pa{ਉੱਤੇ} \pa{ਹਫਦੇ} \pa{ਹਨ।}
\end{quote}

\begin{quote}
\emph{dādā jī pauṛī utte haphde han.}
\end{quote}

\begin{quote}
Grandfather pants on the stairs.
\end{quote}

Past (he …): \pa{ਉਹ} \pa{ਪੌੜੀ} \pa{ਉੱਤੇ} \pa{ਹਫਿਆ}, \emph{oh pauṛī utte haphiā}.

\begin{cousinweb}
\textbf{\pa{ਹਫਣਾ}} (\emph{haphṇā}) is inherited from a reconstructed \emph{*hamph}, \textquotedblleft{}pant\textquotedblright{}; CDIAL lists Punjabi \emph{haphṇā}, \textquotedblleft{}to pant\textquotedblright{}, under it (CDIAL 13973).
\end{cousinweb}

\begin{grammarlens}[title={What you've built}]
One more word for this chapter.
\end{grammarlens}

\subsection*{Your turn}
\begin{itemize}
\raggedright
  \item \emph{Say it:} \emph{haphṇā}
  \item \emph{Say it:} \emph{haphṇā}, once more
  \item \emph{Say it:} \emph{dādā jī pauṛī utte haphde han.} — Grandfather pants on the stairs.
  \item \emph{Recall:} say the Punjabi for an arm, then the Punjabi for a shoulder, then say \emph{haphṇā} again
  \item \emph{Write it:} copy \textbf{\pa{ਹਫਣਾ}} once, checking the whole word
\end{itemize}

\begin{tcolorbox}[breakable,colback=teal!4,colframe=teal!35!black,title={Before you move on}]
What does \pa{ਹਫਣਾ} mean? (\textquotedblleft{}To pant\textquotedblright{}.) Say it once more.
\end{tcolorbox}
\section[\texorpdfstring{\pa{ਪਸੀਨਾ} (pasīnā)}{ਪਸੀਨਾ (pasīnā)}]{\pa{ਪਸੀਨਾ} (pasīnā) — sweat}

\label{lesson:PA-C324-pasina}



\begin{quote}
Before the new one: say the Punjabi for a shoulder, then the Punjabi for to pant.
\end{quote}

\subsection*{You'll want to know: \pa{ਪਸੀਨਾ}}
\textbf{\pa{ਪਸੀਨਾ}} — \emph{pasīnā} — \textquotedblleft{}sweat\textquotedblright{}.

\begin{quote}
\pa{ਗਰਮੀ} \pa{ਵਿੱਚ} \pa{ਮੈਨੂੰ} \pa{ਬਹੁਤ} \pa{ਪਸੀਨਾ} \pa{ਆਉਂਦਾ} \pa{ਹੈ।}
\end{quote}

\begin{quote}
\emph{garmī vich mainū̃ bahut pasīnā āundā hai.}
\end{quote}

\begin{quote}
In the heat I sweat a lot.
\end{quote}

A masculine noun.

\begin{cousinweb}
\textbf{\pa{ਪਸੀਨਾ}} (\emph{pasīnā}) is inherited from a reconstructed \emph{*prasvidana}, \textquotedblleft{}sweating\textquotedblright{}; CDIAL lists Punjabi \emph{pasīṇā}, \textquotedblleft{}sweat\textquotedblright{}, under it (CDIAL 8895).
\end{cousinweb}

\begin{grammarlens}[title={What you've built}]
One more word for this chapter.
\end{grammarlens}

\subsection*{Your turn}
\begin{itemize}
\raggedright
  \item \emph{Say it:} \emph{pasīnā}
  \item \emph{Say it:} \emph{pasīnā}, once more
  \item \emph{Say it:} \emph{garmī vich mainū̃ bahut pasīnā āundā hai.} — In the heat I sweat a lot.
  \item \emph{Recall:} say the Punjabi for a shoulder, then the Punjabi for to pant, then say \emph{pasīnā} again
  \item \emph{Write it:} copy \textbf{\pa{ਪਸੀਨਾ}} once, checking the whole word
\end{itemize}

\begin{tcolorbox}[breakable,colback=teal!4,colframe=teal!35!black,title={Before you move on}]
What does \pa{ਪਸੀਨਾ} mean? (\textquotedblleft{}Sweat\textquotedblright{}.) Say it once more.
\end{tcolorbox}
\section[\texorpdfstring{\pa{ਚੱਕਰ} (chakkar)}{ਚੱਕਰ (chakkar)}]{\pa{ਚੱਕਰ} (chakkar) — a dizzy spell}

\label{lesson:PA-C324-chakkar}



\begin{quote}
Before the new one: say the Punjabi for to pant, then the Punjabi for sweat.
\end{quote}

\subsection*{You'll want to know: \pa{ਚੱਕਰ}}
\textbf{\pa{ਚੱਕਰ}} — \emph{chakkar} — \textquotedblleft{}a dizzy spell\textquotedblright{}.

\begin{quote}
\pa{ਕੱਲ੍ਹ} \pa{ਗਰਮੀ} \pa{ਵਿੱਚ} \pa{ਮੈਨੂੰ} \pa{ਚੱਕਰ} \pa{ਆਇਆ।}
\end{quote}

\begin{quote}
\emph{kall garmī vich mainū̃ chakkar āiā.}
\end{quote}

\begin{quote}
Yesterday I felt dizzy in the heat.
\end{quote}

A masculine noun.

\begin{cousinweb}
\textbf{\pa{ਚੱਕਰ}} (\emph{chakkar}) is inherited from Sanskrit \emph{cakra}, \textquotedblleft{}a wheel, a circle\textquotedblright{}; CDIAL lists Punjabi \emph{cakkar}, \textquotedblleft{}circuit, ring\textquotedblright{}, under it (CDIAL 4538). The sense of a dizzy spell is not given there.
\end{cousinweb}

\begin{grammarlens}[title={What you've built}]
One more word for this chapter.
\end{grammarlens}

\subsection*{Your turn}
\begin{itemize}
\raggedright
  \item \emph{Say it:} \emph{chakkar}
  \item \emph{Say it:} \emph{chakkar}, once more
  \item \emph{Say it:} \emph{kall garmī vich mainū̃ chakkar āiā.} — Yesterday I felt dizzy in the heat.
  \item \emph{Recall:} say the Punjabi for to pant, then the Punjabi for sweat, then say \emph{chakkar} again
  \item \emph{Write it:} copy \textbf{\pa{ਚੱਕਰ}} once, checking the whole word
\end{itemize}

\begin{tcolorbox}[breakable,colback=teal!4,colframe=teal!35!black,title={Before you move on}]
What does \pa{ਚੱਕਰ} mean? (\textquotedblleft{}A dizzy spell\textquotedblright{}.) Say it once more.
\end{tcolorbox}
